% =========================================================================
% SUPPLEMENTARY MATERIAL: Reproducibility Manifest, Artifact Map,
% Data Model Specifications, and Verifier Issue Taxonomy
%
% Accompanies: "A Verifier-Gated Neuro-Symbolic Agent Creates Isomorphic Math
%               Variants and Step Scaffolds"
% Target Journal: PeerJ Computer Science
% =========================================================================

\documentclass[11pt,letterpaper]{article}
\usepackage[utf8]{inputenc}
\usepackage[T1]{fontenc}
\usepackage{lmodern}
\usepackage{amsmath,amssymb}
\usepackage{geometry}
\geometry{margin=1in}
\usepackage{booktabs,longtable,array,tabularx}
\usepackage{xcolor}
\usepackage{hyperref}
\hypersetup{
    colorlinks=true,
    linkcolor=blue,
    citecolor=blue,
    urlcolor=blue,
    pdftitle={Supplementary Material: Reproducibility Manifest and Artifact Map}
}

\setlength{\parindent}{0pt}
\setlength{\parskip}{6pt plus 2pt minus 1pt}

\title{\textbf{Supplementary Material: Reproducibility Manifest, Data Models, and Verifier Issue Taxonomy}}
\author{\textbf{Antigravity Research Team}}
\date{\today}

\begin{document}

\maketitle

This supplementary document accompanies the manuscript \textbf{``A Verifier-Gated Neuro-Symbolic Agent Creates Isomorphic Math Variants and Step Scaffolds''}. It provides full technical specifications, dataset provenance, random seeds, serialized data schemas, LLM execution configurations, verifier issue taxonomy, and post-hoc answer-audit protocols referenced in the main text.

\tableofcontents
\newpage

% -------------------------------------------------------------------------
% SECTION 1: EXPERIMENT IDENTIFIERS AND SAMPLING SEEDS
% -------------------------------------------------------------------------
\section{Experiment Identifiers and Sampling Seeds}\label{sec:experiment-identifiers}

To ensure complete reproducibility of the empirical evaluation reported in Section 4, all experimental trials were executed against a fixed sample derived from deterministic pseudo-random shuffling.

\begin{itemize}
    \item \textbf{Experiment Identifier:} \texttt{tri\_dataset\_200\_seed\_20260803}
    \item \textbf{Experiment Directory:} \texttt{experiments/tri\_dataset\_200\_seed\_20260803/}
    \item \textbf{Primary Sampling Seed:} \texttt{20260803}
    \item \textbf{Dataset-Specific Shuffling Seeds:}
    \begin{itemize}
        \item \textbf{GSM8K:} \texttt{20260803}
        \item \textbf{MATH Algebra:} \texttt{20260804}
        \item \textbf{PRM800K:} \texttt{20260805}
    \end{itemize}
    \item \textbf{Experimental Design:} 600 source items (200 items/dataset $\times$ 3 datasets) $\times$ 2 generation modes (\texttt{isomorphic}, \texttt{scaffold}) $\times$ 2 LLM models (\texttt{gpt-4o-mini}, \texttt{deepseek-v3}) = 2,400 total experimental trials.
\end{itemize}

% -------------------------------------------------------------------------
% SECTION 2: DATASET PROVENANCE AND SHA-256 MANIFEST
% -------------------------------------------------------------------------
\section{Dataset Provenance and SHA-256 Manifest}\label{sec:dataset-provenance}

Table~\ref{tab:dataset-provenance} details the upstream source endpoints, local file paths, and cryptographic SHA-256 hashes for all evaluation datasets.

\begin{table}[htbp]
\centering
\small
\caption{Dataset provenance, local storage paths, and cryptographic SHA-256 manifests.}
\label{tab:dataset-provenance}
\begin{tabular}{lllll}
\toprule
Dataset & Upstream Source & Split / Config & Local Path & SHA-256 Checksum \\
\midrule
GSM8K & Hugging Face \texttt{openai/gsm8k} & \texttt{main}, test split & \texttt{data/raw/gsm8k\_test.jsonl} & \texttt{752adc99f...abafbac} \\
MATH & Hugging Face \texttt{EleutherAI/hendrycks\_math} & \texttt{algebra}, test split & \texttt{data/raw/math\_algebra\_test.jsonl} & \texttt{81a4218fc...89fcc03} \\
PRM800K & OpenAI \texttt{openai/prm800k} & \texttt{phase2\_test} & \texttt{data/raw/prm800k\_phase2\_test.jsonl} & \texttt{6b172efa8...d98aaad} \\
\bottomrule
\end{tabular}
\end{table}

\textbf{Direct Download Endpoints:}
\begin{itemize}
    \item \textbf{GSM8K:} \url{https://raw.githubusercontent.com/openai/grade-school-math/master/grade_school_math/data/test.jsonl}
    \item \textbf{MATH Algebra:} \url{https://people.eecs.berkeley.edu/~hendrycks/MATH.tar.gz}
    \item \textbf{PRM800K:} \url{https://media.githubusercontent.com/media/openai/prm800k/main/prm800k/data/phase2_test.jsonl}
\end{itemize}

% -------------------------------------------------------------------------
% SECTION 3: LLM PROVIDER AND RUNTIME SERVICE CONFIGURATION
% -------------------------------------------------------------------------
\section{LLM Provider and Service Configurations}\label{sec:service-configuration}

Table~\ref{tab:model-service-configuration} summarizes the LLM service endpoints, sampling parameters, and context limits used during formalization and candidate generation.

\begin{table}[htbp]
\centering
\small
\caption{LLM service endpoints and hyperparameter configurations.}
\label{tab:model-service-configuration}
\begin{tabular}{lllrr}
\toprule
Pipeline Stage & Reported Model Identifier & Default Base URL & Temp. & Max Tokens \\
\midrule
Formalization (Decomposer) & \texttt{gpt-4o-mini} & \url{https://api.openai.com/v1} & 0.2 & 2,500 \\
GPT Candidate Generation & \texttt{gpt-4o-mini} & \url{https://api.openai.com/v1} & 0.2 & 1,800 \\
DeepSeek Candidate Gen. & \texttt{deepseek-ai/DeepSeek-V3} & \url{https://api.siliconflow.cn/v1} & 0.2 & 1,800 \\
\bottomrule
\end{tabular}
\end{table}

\textbf{Additional Runtime Execution Parameters:}
\begin{itemize}
    \item \textbf{Concurrency Limit:} 4 parallel worker threads.
    \item \textbf{Maximum Agent Retries ($k$):} 2 retry attempts per trial (up to 3 total attempts including initial generation).
    \item \textbf{JSON Schema Repair Allowance:} Up to 2 automatic repair requests per API call if JSON parsing fails.
    \item \textbf{API Timeout:} 60 seconds per HTTP request.
    \item \textbf{Credential Isolation:} API keys are retrieved from environment variables (\texttt{OPENAI\_API\_KEY}, \texttt{SILICONFLOW\_API\_KEY}) or git-ignored \texttt{config.json} files, and are scrubbed from log files and public artifacts.
\end{itemize}

% -------------------------------------------------------------------------
% SECTION 4: SERIALIZED DATA MODELS AND CONTROLLER STATES
% -------------------------------------------------------------------------
\section{Serialized Data Models and Controller States}\label{sec:data-models}

The controller architecture formalizes problem items, steps, candidate responses, and verifier outputs using immutable data structures.

\subsection{Controller Terminal States}
The \texttt{MathAgent} controller outputs one of two terminal states:
\begin{enumerate}
    \item \textbf{\texttt{DELIVERED(candidate, verification\_result, attempt\_history)}}: Emitted when a generated candidate passes all SymPy verifier checks ($CVI \ge 0.60$, $SFS \ge 0.60$, real solvable roots, satisfied domain constraints).
    \item \textbf{\texttt{ABSTAINED(attempt\_history, failure\_reasons)}}: Emitted when all $k$ attempts fail verification, triggering fail-safe abstention to prevent invalid items from reaching students.
\end{enumerate}

\subsection{Serialized Data Schema}
\begin{itemize}
    \item \textbf{\texttt{Problem}}: \texttt{\{problem\_id, text, variables, equations, steps, domain\_constraints\}}
    \item \textbf{\texttt{Step}}: \texttt{\{step\_id, name, description, equations, dependencies, concepts\}}
    \item \textbf{\texttt{Candidate}}: \texttt{\{question, variables, equations, answer, reasoning\}}
    \item \textbf{\texttt{VerificationResult}}: \texttt{\{accepted, solutions, cvi, isomorphism\_score, issues\}}
    \item \textbf{\texttt{Attempt}}: \texttt{\{number, candidate, verification, reflection, error\}}
\end{itemize}

% -------------------------------------------------------------------------
% SECTION 5: VERIFIER FAILURE ISSUE TAXONOMY AND REFLECTION PROMPTS
% -------------------------------------------------------------------------
\section{Verifier Failure Issue Taxonomy and Reflection Mappings}\label{sec:issue-taxonomy}

When candidate verification fails, the \texttt{SymbolicVerifier} categorizes the failure into an issue taxonomy code. The \texttt{ReflectionEngine} maps these codes into prompt constraints for the next attempt (Table~\ref{tab:issue-taxonomy}).

\begin{table}[htbp]
\centering
\small
\caption{Verifier issue taxonomy codes and corresponding reflection prompt actions.}
\label{tab:issue-taxonomy}
\begin{tabular}{lll}
\toprule
Issue Code & Failure Description & Reflection Prompt Action \\
\midrule
\texttt{PARSE\_ERROR} & Equation or JSON schema cannot be parsed & Enforce strict SymPy syntax and declared variables. \\
\texttt{NO\_SOLUTION\_ERROR} & Equation system has no finite solution & Add numerical givens to uniquely determine variables. \\
\texttt{SOLVER\_ERROR} & SymPy internal solver exception or timeout & Simplify algebraic structure and eliminate circularity. \\
\texttt{NON\_REAL\_SOLUTION\_ERROR} & Solutions contain complex imaginary parts & Ensure discriminant $\Delta \ge 0$ for real-valued roots. \\
\texttt{DOMAIN\_VIOLATION\_ERROR} & Solutions violate active constraints ($x \le 0$) & Regenerate coefficients so roots satisfy predicates. \\
\texttt{UGLY\_SOLUTIONS\_ERROR} & Clean-Value Index ($CVI < 0.60$) is low & Restructure coefficients for small integers/fractions. \\
\texttt{NON\_ISOMORPHIC\_ERROR} & Structural similarity ($SFS < 0.60$) is low & Preserve variable count, degree, and operation profile. \\
\texttt{TARGET\_STEP\_MISMATCH} & Scaffold targets wrong step ID & Align step equations strictly with \texttt{target\_step\_id}. \\
\texttt{GENERATION\_ERROR} & API network error or JSON failure & Re-prompt LLM with strict JSON schema instructions. \\
\bottomrule
\end{tabular}
\end{table}

% -------------------------------------------------------------------------
% SECTION 6: POST-HOC ANSWER-AUDIT IMPLEMENTATION
% -------------------------------------------------------------------------
\section{Post-Hoc Answer-Audit Implementation}\label{sec:answer-audit}

To evaluate the mathematical consistency between the natural-language text generated by LLMs and the SymPy-verified solution, a post-hoc audit was implemented (\texttt{scripts/audit\_paper\_results.py}).

\subsection{Audit Classification Categories}
Each delivered candidate item was audited and assigned to one of three categories:
\begin{enumerate}
    \item \textbf{Matched (Answer Agreement):} The candidate's natural-language answer text (e.g., ``$x = 18$'') mechanically parses into a numeric value or tuple that matches the SymPy solver's exact solution set.
    \item \textbf{Confirmed Contradiction:} The parsed natural-language answer text explicitly contradicts every solution computed by SymPy (e.g., text states ``$x = 25$'' but SymPy solver computed ``$x = 18$'').
    \item \textbf{Mechanically Unverified:} The candidate's text contains prose, complex units, or non-standard formatting that cannot be unambiguously parsed by regular expressions, retained in the denominator for conservative accuracy bounds.
\end{enumerate}

\subsection{Statistical Confidence Intervals}
Proportions for contradiction rates and match rates are reported with \textbf{Wilson score 95\% confidence intervals}:
\begin{equation}
    CI_{95\%} = \frac{\hat{p} + \frac{z^2}{2n} \pm z \sqrt{\frac{\hat{p}(1-\hat{p})}{n} + \frac{z^2}{4n^2}}}{1 + \frac{z^2}{n}}
\end{equation}
where $z = 1.96$, $\hat{p}$ is the sample proportion, and $n$ is the total delivered item count.

% -------------------------------------------------------------------------
% SECTION 7: ARTIFACT DIRECTORY MAP AND REGENERATION COMMANDS
% -------------------------------------------------------------------------
\section{Artifact Directory Map and Regeneration Commands}\label{sec:artifact-map}

\subsection{Artifact Directory Structure}
\begin{table}[htbp]
\centering
\small
\caption{Artifact directory map for experimental outputs.}
\label{tab:artifact-directory-map}
\begin{tabular}{ll}
\toprule
Relative Location & Description of Retained Contents \\
\midrule
\texttt{data/benchmark/} & Fixed 200-item paper benchmark JSONL files (\texttt{gsm8k\_200.jsonl}, etc.) \\
\texttt{data/normalized/} & Full normalized test splits for GSM8K, MATH Algebra, and PRM800K \\
\texttt{web/} & Standalone FastAPI server, HTML template, CSS styling, and JS visual renderer \\
\texttt{examples/demo\_run.py} & Dual-mode (Offline + Online API) Python demonstration script \\
\texttt{tests/} & Automated unit test suite (\texttt{test\_agent.py}, \texttt{test\_datasets.py}) \\
\texttt{scripts/run\_large\_experiment.py} & Benchmark execution pipeline for 2,400 experimental trials \\
\texttt{scripts/analyze\_experiment\_for\_paper.py} & Statistical analysis, McNemar tests, CVI/SFS distribution tables \\
\texttt{scripts/audit\_paper\_results.py} & Post-hoc answer-consistency and evidence-chain audit script \\
\bottomrule
\end{tabular}
\end{table}

\subsection{Commands for Full Experiment Regeneration}
To execute the full verification and analysis pipeline from a fresh installation:

\begin{verbatim}
# 1. Install package in editable mode
pip install -e .

# 2. Run unit test suite
python -m unittest discover -s tests -v

# 3. Run interactive Web UI server
python -m uvicorn web.server:app --host 0.0.0.0 --port 8000

# 4. Run statistical analysis and paper table generation
python scripts/analyze_experiment_for_paper.py

# 5. Run post-hoc answer consistency audit
python scripts/audit_paper_results.py
\end{verbatim}

\end{document}
